% reconfiguration.tex -- Methodology: reconfiguration (detach and attach). Draft 4 Oct 2026.
% Code: dissipative_planner/dissipative_node.py (_detach_ocp, _attach_ocp, _capture_weld_rho,
% partner hand-off), dissipative_planner/approach.py, drone_magnet (magnet_attachment_manager,
% elrs_mux, handover_policy, attach_target_publisher), bringup/bringup/graph_attach.py.
% Evidence: CURRENT_STATE 4.1, 4.2, 4.11; notes 2026-09-24_detach_land_fix, 2026-09-25_attach_ocp,
% 2026-09-27_m1_rejoin, 2026-09-27_m1_full, 2026-09-27_m2_handover.

\subsection{Reconfiguration by OCP resize}
\label{sec:reconfig}

A reconfiguration changes the number of UAVs carrying the payload in the air. In the final design both
directions are handled the same way: the planner switches to a solver for the new team size and carries
on. This subsection describes the detach, the attach, and the design each replaced.

\subsubsection{Detach}

On DETACH the departing UAV's magnet releases the rod at the payload end, the UAV flies away with its
rod, and the planner continues with the survivors. Two approaches were built and compared.

\paragraph{OCP resize.} The planner keeps flying the payload and re-solves for the smaller fleet, which
two properties of the formulation make possible.
\begin{itemize}
    \item The attachment points $\boldsymbol{\rho}_i$ and rod lengths $l_i$ are runtime parameters of
    the model, not constants compiled into the solver. Only the fleet size fixes the problem
    dimensions, so one solver per fleet size, built in advance, covers any attachment layout.
    \item The survivors keep their own attachment points. The rods stay where they were on the ring,
    so three survivors of a four-UAV ring sit, for example, at 0, 90 and 180$^\circ$, not at an even
    0, 120 and 240$^\circ$. Handing the planner an even ring instead makes its moment balance wrong by
    construction.
    \redo{the hand-back to an even ring after a detach from four UAVs (the survivors at 0, 90 and 180$^\circ$, the planner told 0, 120 and 240$^\circ$): payload tilt.}
% --- original text, flown before the rig twin; kept to refill this place:
%     0, 120 and 240$^\circ$. Handing the planner an even ring instead makes its moment balance wrong by
%     construction; in an early test this tipped the payload by 62$^\circ$.
\end{itemize}
At the detach the planner swaps to the solver for $N-1$ UAVs, sets the survivors' attachment points and
rod lengths, and warm-starts the new problem from the measured state. The tensions and rod directions are
free at node~0, so the first solve redistributes the load, with no separate hand-off phase. A detach is
refused if it would leave fewer than three UAVs, and during LAND.

\paragraph{Network hand-off.} The alternative, closer to the decentralised literature
\cite{quan2025swarm}, hands the payload to the dissipative network of Section~\ref{sec:dissipative} on the
detach; the network holds it while it redistributes the load between the survivors, then hands back to a
planner built for $N-1$. Section~\ref{sec:sim_detach} compares the two; the OCP resize is the detach used
in the rest of this work.
% --- original text, flown before the rig twin; kept to refill this place:
% then hands back to a planner built for $N-1$. In simulation the OCP resize won this comparison
% clearly (Section~\ref{sec:sim_detach}), and it is the detach used in the rest of this work.

% Departed landing (pressed to its rest, disarmed alone): twin 9 Oct, card 2026-10-09_departed_disarm_alone;
% the default once the LAND re-fly's reviewer passes it; rig W4/W5 on 14 Oct.
\paragraph{The departing UAV.} After the release, switched on the rig by the same DETACH command on the
UAV's magnet channel, the departing UAV steps 0.5\,m out radially at its height, clear of the ring and the
survivors, and lands on its own while the survivors carry on. Its height reference is pressed below its
own resting height until it gets there, and it is then disarmed alone, without waiting for the fleet
(Section~\ref{sec:res_detach_rig}).

\paragraph{Layout.} Whether the payload can still hang level after a detach depends on the geometry, not
on the controller: its centre of mass must stay inside the polygon of the surviving attachment points.
Removing one UAV from an even ring of $N$ leaves a largest angular gap of $4\pi/N$, which keeps the centre
inside only for $N \ge 5$; $N = 4$ sits exactly on the boundary. The four-UAV flights therefore use the
plates at 30, 90, 150 and 270$^\circ$, and detaching the UAV at 90$^\circ$ leaves an even triangle under
which the ring hangs level. With an uneven layout the planner's nominal tensions come from a moment
balance at the actual attachment points, so they are unequal and the payload's reference stays level.

\paragraph{Landing after a detach.} Touchdown is judged on the survivors only, so a departed UAV still
descending elsewhere does not hold their references down; a departed UAV is judged down at its own
resting height, and the fleet is declared landed only when it is down too. The cable feedforward is
switched off once the payload rests, because its sideways component pushes a grounded survivor across
the floor.
% --- original text, flown before the rig twin; kept to refill this place:
% declared landed only when they are down too. The cable feedforward is switched off once the payload
% rests, because its sideways component pushed a grounded survivor across the floor until it tipped
% in an early version.

\subsubsection{Attach}

On ATTACH a free UAV (the newcomer) approaches its attachment point, its magnet welds to a plate on the
ring, and the planner folds it into the team. The approach controller, the approach UAV and its magnet
are the collaborating student's; this work takes over at the weld.
% TODO: cite the collaborator's thesis or report once available.

\paragraph{Sequence.}
\begin{enumerate}
    \item \textbf{Target hold.} ATTACH freezes the payload's horizontal reference, or keeps it on its
    trajectory when the payload is moving, shifted onto the measured payload position, so the newcomer
    approaches where the payload actually is.
    \item \textbf{Approach.} The newcomer flies to a point above its plate and descends onto it, flown by
    the collaborator's approach controller or by its own tracker on an approach profile from this work.
    The payload may keep moving, at up to 0.2\,m/s.
    \redo{a weld onto a payload moving at 0.4\,m/s against 0.2\,m/s (the limit on the payload's speed during the approach).}
% --- original text, flown before the rig twin; kept to refill this place:
%     approach profile from this work. The payload may keep moving, at up to 0.2\,m/s; at 0.4\,m/s the
%     weld flipped the payload in simulation.
    \item \textbf{Weld.} A magnet manager welds when the magnet tip is within 8\,cm of the plate point,
    computed from the payload's measured pose, and the relative speed is under 0.05\,m/s, both for 0.15\,s.
    \item \textbf{Authority hand-over.} Until the weld, a multiplexer forwards the approach controller's
    commands to the newcomer's radio; after it, the multiplexer switches to this work's tracker, only on a
    command that is armed and above idle throttle, and latches. An abort latches the mux to a disarmed idle
    command until it is restarted.
    \redo{the authority hand-over at the weld with the mux switching on an armed but idle stream, against the armed-and-above-idle rule: whether the newcomer's motors cut at the weld.}
% --- original text, flown before the rig twin; kept to refill this place:
%     tracker, but only on a command that is armed and above idle throttle. A mux that handed authority
%     to an armed but idle stream cut the newcomer's motors at every weld in early tests. The switch is
%     latched. An abort latches the mux to a disarmed idle command until it is restarted.
    \item \textbf{Fold-in.} The planner records the weld point in the payload frame and switches to
    the solver for $N+1$ UAVs, with the newcomer's tension starting at 0.1\,N and every UAV's tension
    blended to the balanced solution for the new geometry.
    \redo{the fold-in with the newcomer's tension starting at 1\,N against 0.1\,N: payload tilt at the weld.}
% --- original text, flown before the rig twin; kept to refill this place:
%     blended to the balanced solution for the new geometry. A newcomer started at 1\,N of tension
%     pulled the payload with a lever about the incumbents' cable apex and tilted it at the weld.
    \item \textbf{Resume.} The trajectory continues with $N+1$ UAVs.
\end{enumerate}

\paragraph{Replaced design: network attach.} Until mid-September the attach handed the payload to the
dissipative network at the weld: the newcomer was welded as a central lifter and handed out to its ring
slot over 8--12\,s, with the incumbents' tensions blended over the same ramp
(compared in Section~\ref{sec:sim_attach}).
% --- original text, flown before the rig twin; kept to refill this place:
% its ring slot over 8--12\,s, with the incumbents' tensions blended over the same ramp. This flew the
% demonstration of Figure~\ref{fig:attach} on the earlier, lighter ring. On the 0.86\,kg ring the
% network capsized the payload within about 8\,s of every weld in simulation, while the OCP resize held
% the same weld at about 10$^\circ$ (Section~\ref{sec:sim_attach}).

\paragraph{The rejoin mission.} The full mission of Figure~\ref{fig:mission_profile} is assembled with the
collaborator's mission planner. Four UAVs carry the ring around the orbit; one detaches by OCP resize, flies
the collaborator's mission (picking up an object with its magnet and dropping it into the ring), is handed
back at its attach-ready point, approaches, welds and is folded back in by OCP resize, and the fleet lands.
In a ground-start variant the collaborator's controller joins four UAVs to the ring on the floor and hands
them to this work's planner, which lifts and carries it.
